\chapter{How Are You?}
\label{ch:responding}

The everyday exchange after hello:

\begin{center}
\mr{तुम्ही कसे आहात?} \quad(\emph{tumh\=\i\ kase \=ah\=at} --- ``How are you?'')\\
\mr{मी बरा आहे, धन्यवाद.} \quad(\emph{m\=\i\ bar\=a \=ahe, dhanyav\=ad} --- ``I'm well, thank you.'')
\end{center}

\emph{Practice lessons: \texttt{lessons/MR-C03-*.md}.}

\section{\mr{कसा} / \mr{कशी} / \mr{कसं} \emph{(kas\=a / ka\'{s}\=\i\ / kas\d{m})} --- ``how''}
\label{lesson:kasa}

\begin{cousinweb}
\mr{कसा} (\emph{kas\=a}, ``how'') is another \emph{k-} question, from PIE
\emph{*k\ensuremath{^w}o-} --- the root behind English \textbf{how} (a \emph{wh-}
word that lost its \emph{w}) and Latin \emph{quo-modo} (``in what way,'' $\to$
Spanish \emph{c\'{o}mo}).
\end{cousinweb}

\begin{grammarlens}
\textbf{``How'' agrees with whom you ask.} \mr{कसा} (\emph{kas\=a}, to a man) /
\mr{कशी} (\emph{ka\'{s}\=\i}, to a woman) / \mr{कसं} (\emph{kas\d{m}}, neuter). So
``how are you?'' shifts shape with the person addressed.
\end{grammarlens}

\section{\mr{तुम्ही कसे आहात?} --- ``how are you?''}
\label{lesson:tumhi-kase-aahat}

\mr{तुम्ही} (\emph{tumh\=\i}, resp.\ ``you'') $+$ \mr{कसे} (\emph{kase}, ``how,''
plural-agreeing) $+$ \mr{आहात} (\emph{\=ah\=at}, ``are,'' resp.) $=$ ``you how
are?'', verb last. Because \emph{tumh\=\i} is grammatically plural, ``how''
appears as \emph{kase} and the copula as \emph{\=ah\=at}. To a male friend:
\mr{तू कसा आहेस?} (\emph{t\=u kas\=a \=ahes?}); to a female friend: \mr{तू कशी आहेस?}
(\emph{t\=u ka\'{s}\=\i\ \=ahes?}).

\section{\mr{मी} \emph{(m\=\i)} --- ``I''}
\label{lesson:mi}

\begin{cousinweb}
\mr{मी} (\emph{m\=\i}, ``I'') descends from Sanskrit \mr{अहम्} (\emph{aham}),
worn down to \emph{m\=\i}. That \emph{aham} is PIE
\emph{*h\textsubscript{1}e\'{g}(H)om} --- the ancestor of Latin \textbf{ego} and
English \textbf{I}. Note its family tie to last chapter's \emph{m\=ajh\d{m}}
(``my''): the same first-person \emph{m-} thread, as English pairs \emph{I} with
\emph{me/my}.
\end{cousinweb}

\begin{grammarlens}
\textbf{``I am'' is \emph{m\=\i\ \=ahe}.} And ``I'm fine'' reaches back to
Chapter~1's \emph{bar\d{m}} in its gendered form: \emph{m\=\i\ bar\=a \=ahe}
(a man) / \emph{m\=\i\ bar\=\i\ \=ahe} (a woman).
\end{grammarlens}

\section{\mr{मी बरा आहे} --- ``I'm well, thank you''}
\label{lesson:mi-bara-aahe}

\mr{मी} (I) $+$ \mr{बरा} (\emph{bar\=a}, ``well'') $+$ \mr{आहे} (am) $=$ ``I am
well.'' Chapter~1's \mr{बरं} (\emph{bar\d{m}}) was the bare \textbf{neuter}
``okay''; describing a person it takes their gender: \emph{bar\=a} (m.) /
\emph{bar\=\i} (f.). Add thanks: \emph{m\=\i\ bar\=a \=ahe, dhanyav\=ad}.

\begin{grammarlens}
\textbf{The copula changes with the subject.} \emph{m\=\i\ \=ahe} (I am), \emph{t\=u
\=ahes} (you are), \emph{tumh\=\i\ \=ah\=at} (you are, resp.) --- one verb,
reshaped for each person. Marathi marks \emph{who} on the verb itself.
\end{grammarlens}

\section{\mr{काही हरकत नाही} \emph{(k\=ah\=\i\ harkat n\=ah\=\i)} --- ``no problem''}
\label{lesson:kaahi-harkat-nahi}

\begin{cousinweb}
\mr{काही} (\emph{k\=ah\=\i}, ``any'') $+$ \mr{हरकत} (\emph{harkat}, ``objection'')
$+$ \mr{नाही} (\emph{n\=ah\=\i}, ``not'') $=$ ``[there is] no objection at all''
--- Marathi's easy ``you're welcome / no problem.'' Two histories meet:
\emph{n\=ah\=\i} rides PIE \emph{*ne} (English \textbf{no}, \textbf{not}), while
\emph{harkat} is a \textbf{Perso-Arabic} loan (Arabic \emph{\d{h}arakah}) --- one
of the many Persian and Arabic words Marathi absorbed under the Deccan
sultanates. An Indo-European negative wrapped around an Arabic noun.
\end{cousinweb}

\section{The whole exchange}
\label{lesson:MR-C03-practice}

\begin{center}
\begin{tabular}{@{}ll@{}}
\toprule
Marathi & English \\
\midrule
\mr{तुम्ही कसे आहात?} & How are you? \\
\mr{मी बरा आहे, धन्यवाद.} & I'm well, thank you. \\
\mr{काही हरकत नाही.} & No problem / you're welcome. \\
\bottomrule
\end{tabular}
\end{center}

Roots banked: \emph{kas\=a} (the \emph{k-} questions, gender-agreeing), \emph{m\=\i}
(\emph{aham} $\to$ Latin \emph{ego}, English \emph{I}), \emph{bar\=a/bar\=\i}
(now gendered), \emph{harkat} (Arabic, the Deccan trade-and-court layer). Next
chapter: the farewells --- \emph{punh\=a bhe\d{t}\=u}.
